\section{A conditional exact-calibration benchmark in finite strata}
\label{sec:conditional-standardization}

This is a separate feasible finite-stratum procedure. It changes score
construction and sample use; it is not merely a smaller constant for the
previous fitted-score method. Its purpose is to determine whether useful
borrowing is possible once empirical design imbalance is removed.
It does not establish a high-dimensional or general model-robust RoCE theorem.

\subsection{Conditional target and exact calibration}

Assume the causal identification conditions under which $\mu_t^a(x)=\E_t[Y\mid A=a,X=x]$ is the target potential-outcome mean.
Let $X$ take $L$ known values, and condition on all observed covariates and
treatments, denoted $\mathcal D$. Let $\hat p_{t,x}$ be the full-target empirical
cell proportions. Define
\[
 \mu_{t,X}^a=\sum_x\hat p_{t,x}\mu_t^a(x),\qquad
 \tau_X=\mu_{t,X}^1-\mu_{t,X}^0.
\]
This is the average of conditional treatment effects over the observed target
covariates, not the average of realized counterfactual outcomes.
The population parameter remains $\tau_t=\E_t\{\mu_t^1(X)-\mu_t^0(X)\}$.
For each arm, at least $g_a$ sources share the target conditional outcome means
in every cell. The remaining conditional means are arbitrary in $[0,1]$;
no lower bound on their nonzero deviations is assumed.

If source $j$ has $N_{jxa}$ observed patients in cell $(x,a)$, exact design
calibration gives
\[
 q_j^a(x)=\frac{n_j\hat p_{t,x}}{N_{jxa}},\qquad
 Z_j^a=\sum_x\hat p_{t,x}\bar Y_{jxa}.
\]
For any common cellwise predictor $m^a$, including an estimated predictor,
\[
 \sum_x\hat p_{t,x}m^a(x)+\frac1{n_j}\sum_i
 \ind(A_{ji}=a)q_j^a(X_{ji})\{Y_{ji}-m^a(X_{ji})\}
 =Z_j^a.
\]
Thus the fitted prediction cancels exactly. Given $\mathcal D$, a valid source
has mean $\mu_{t,X}^a$ with no nuisance fitting drift. The weights use designs,
not the evaluated outcomes. This is a saturated instance of balancing and
standardization, not a new weighting principle \citep{sbw2015,arb2018}.

The implementation requires at least two observations in every used source
cell/arm to estimate cell variances unbiasedly. If this fails, it uses the
prespecified target-only conditional band for that repetition. No repetition
is discarded. The fallback decision depends only on designs. One can use
only the original 500 held-out source observations, or all 1,000 observations;
both options are reported. Target outcomes are used for a target conditional
band and the CATE-range certificate. No true outcome means, propensities or
variances enter the procedure.

\subsection{A dispersion bound for independent cell means}

The general source statistic is $Z_j=\sum_h a_{jh}\bar Y_{jh}$, with independent
Bernoulli outcomes within and between cells conditional on $\mathcal D$.
The coefficients can have either sign: an arm mean uses $a_{jh}=\hat p_{t,h}$;
a TATE contrast stacks $+\hat p_{t,h}$ and $-\hat p_{t,h}$.
For $N_{jh}\ge2$ put
\begin{align*}
 \hat v_j&=\sum_h a_{jh}^2
       \frac{\bar Y_{jh}(1-\bar Y_{jh})}{N_{jh}-1},&
 v_{jh}&=\frac{a_{jh}^2}{4N_{jh}},& P_j&=\sum_hv_{jh},\\
 \widehat D&=K^{-1}\sum_j(Z_j-\bar Z)^2
                 -\frac{K-1}{K^2}\sum_j\hat v_j,&
 D&=K^{-1}\sum_j(\E[Z_j\mid\mathcal D]-\E[\bar Z\mid\mathcal D])^2.
\end{align*}
Then $\E(\widehat D\mid\mathcal D)=D$. Define
\begin{align*}
 L_D&=4\max_jP_j/K,\\
 Q_D&=\frac2{K^4}\left[(\sum_jP_j)^2+
       \{(K-1)^2-1\}\sum_jP_j^2+
       (K-1)^2\sum_{jh}\frac{v_{jh}^2}{N_{jh}-1}\right],\\
 c_D&=\frac{2(K-1)}{K^2}\max_{jh}\frac{v_{jh}}{N_{jh}-1}.
\end{align*}
Use the earlier exponential inversion with $L_D,Q_D,c_D$ to get $U_D$ at
failure probability $\alpha_D$.

\begin{proposition}[Conditional cell-mean dispersion]
Under independent bounded binary outcomes given the designs, the preceding
$U_D$ satisfies $\Prb(D\le U_D\mid\mathcal D)\ge1-\alpha_D$.
For $K=1$, set $D=U_D=0$.
\end{proposition}
\begin{proof}
Set $A=I/K-\boldsymbol1\boldsymbol1^T/K^2$. The patient quadratic matrix has
zero diagonal. Within the same source cell its off-diagonal entries are
$A_{jj}a_{jh}^2/\{N_{jh}(N_{jh}-1)\}$; between different cells they are
$A_{jk}a_{jh}a_{k\ell}/(N_{jh}N_{k\ell})$.
The latter also applies to different cells in the same source.
For centered patient outcomes, the linear term has entries
$a_{jh}(A\theta)_j/N_{jh}$, where $\theta_j=\E[Z_j\mid\mathcal D]$.
This follows from constant conditional means within each cell and is valid
when those means differ across cells.

Replace centered Bernoulli variables in the MGF by independent Gaussians with
proxy variance $1/4$, using the zero-diagonal coordinatewise argument.
On cell-constant vectors the scaled quadratic matrix is positive semidefinite;
its nonzero spectrum is that of
$\diag(\sqrt P)A\diag(\sqrt P)$. On each cell's orthogonal complement it
has eigenvalue $-A_{jj}v_{jh}/(N_{jh}-1)$ with multiplicity $N_{jh}-1$.
The Gaussian linear vector is in the positive subspace, has squared norm at
most $(\max_jP_j)D/K$, and the squared Frobenius norm is $Q_D/2$.
Consequently
\[
 \log\E\{e^{-t(\widehat D-D)}\mid\mathcal D\}
 \le \frac{t^2(L_DD+Q_D)}{2(1-c_Dt)},\quad 0<t<1/c_D.
\]
Chernoff's bound and the previously stated quadratic inversion give the result.
\end{proof}

\subsection{One conditional compatibility procedure}

The source center is $T_S=K^{-1}\sum_j(Z_j^1-Z_j^0)$. Its outcome-noise
proxy is known from the design:
\[
 V_S^+=K^{-2}\sum_{jxa}\frac{\hat p_{t,x}^2}{4N_{jxa}},\qquad
 r_S=\sqrt{2V_S^+\log(2/\alpha_S)}.
\]
There is no claim of unconditional independence between the two arms. Here
conditioning on all designs leaves outcomes in disjoint treatment cells
independent; the common target-covariate uncertainty is handled separately.

The arm dispersion certificate is
\[
 B_A=\sqrt{(K-g_1)U_1/g_1}+\sqrt{(K-g_0)U_0/g_0}.
\]
If $g_\cap=g_1+g_0-K>0$, the TATE certificate is
$B_C=\sqrt{(K-g_\cap)U_\tau/g_\cap}$. The adaptive option allocates half
its dispersion error to the two arm certificates and half to the TATE
certificate, then uses $\min(B_A,B_C)$. This data-dependent minimum is
protected by the simultaneous events. Without a guaranteed intersection,
it uses only $B_A$ and reallocates the dispersion allowance by a fixed rule.
Thus arm-specific valid sets are preserved.

The source conditional band is $J_S=[T_S\pm(r_S+B)]$.
The target conditional band $J_T$ estimates $\tau_X$ with all target outcomes,
using unsmoothed cell means when an arm is observed and .5 otherwise. Empty
arms contribute an explicit worst-case bias allowance. A weighted Hoeffding
bound provides its conditional outcome radius.
Intersect $J_S,J_T$ when possible and use the shorter conditional band when
the intersection is empty. If source cells are unavailable, use $J_T$.

To pass from $\tau_X$ to $\tau_t$, let $R^+$ be a simultaneous upper bound
on the true CATE range from target cellwise binomial intervals. With target
sample size $n_t$, add
\[
 r_X=R^+\sqrt{\log(2/\alpha_X)/(2n_t)}
\]
to both conditional endpoints and clip to $[-1,1]$.
The procedure handles the common target composition uncertainty once.
Its fixed failure allocation is
\[
 \underbrace{.005}_{\text{range}}+
 \underbrace{.01}_{\text{composition}}+
 \underbrace{.01}_{\text{source outcomes}}+
 \underbrace{.01}_{\text{target outcomes}}+
 \underbrace{.015}_{\text{dispersion}}=.05.
\]
On the simultaneous events the conditional intersection contains $\tau_X$,
and enlargement contains $\tau_t$. No independence between those events is
needed. The design-dependent unavailable-source branch uses the same target
outcome, range and composition events.

The comparison target-only interval uses the same unsmoothed stratified point
and range certificate, with .005 range error and .0225 each for outcome and
composition error. It therefore has its own .95 finite-sample guarantee.

\subsection{A point rule with an explicit conditional property}

Keep the midpoint as a recorded summary, and optionally project the target
point estimate onto the selected conditional compatibility band before
population enlargement. On the event $\tau_X\in J_S$, this projection cannot
increase absolute error relative to the unprojected target point for $\tau_X$.
Indeed the target point lies in its own band; projection onto a nonempty
intersection equals projection onto $J_S$. If an empty intersection selects
$J_S$, the same contraction applies; if it selects $J_T$, the target point
is unchanged. Clipping to the known effect range preserves this property.
This is not a universal MSE guarantee for the population parameter $\tau_t$;
target composition error and certificate failures still matter. Both point
rules are evaluated with new seeds, without choosing the better rule separately
in each scenario.

\subsection{Scope, rates, and the next extension}

This benchmark assumes finite strata and conditional-mean transport for the
guaranteed valid sites. It does not transfer automatically to cases where only
a population-average source moment is valid. For fixed $L$, uniformly adequate
cell counts, $g_a/K$ bounded away from zero, and zero conditional dispersion,
$P_j=O(n_s^{-1})$, $Q_D=O(n_s^{-2}K^{-1})$, and
$c_D=O(n_s^{-2}K^{-1})$. The source bias allowance has order
$n_s^{-1/2}K^{-1/4}$, while source outcome noise is
$O(n_s^{-1/2}K^{-1/2})$. Target composition remains of order $n_t^{-1/2}$
when the effect is heterogeneous. Uniformly adequate random cell counts
require a joint sample-size/source-count condition; $K\to\infty$ with fixed
per-site sample size alone does not guarantee it. Coverage remains protected
by the target fallback, but useful source efficiency need not.

Exact balance removes the fitted remainder in this saturated experiment.
The relevant next question is how approximate balance and regularized outcome
fits change that remainder and its noise correction uniformly in $K$.
\citet{arb2018} gives a high-dimensional linear-outcome foundation; a recent
nonlinear residual-balancing preprint studies derivative and curvature control
\citep{meza2025}. Those results require their own assumptions and have not been
proved to deliver the present multi-source conditional procedure. Balancing
weights themselves are established methodology; the open contribution is
uniform inference with unknown arm-specific valid sources and weak deviations.
